La implementación de un sistema de información con capacidades predictivas y comparativas facilitará el análisis del desempeño académico de los estudiantes y aportará una visión más clara sobre su proceso de formación~\cite{granados2025modelado}. Este proyecto busca integrar los resultados de las pruebas Saber 11 y Saber Pro en una única plataforma de análisis que permita comprender cómo evolucionan las competencias de los estudiantes desde su ingreso a la universidad hasta la finalización de sus estudios. A través de esta herramienta, la Universidad de San Buenaventura podrá identificar patrones de mejora o posibles áreas de estancamiento en el desarrollo de competencias, así como reconocer los factores que tienen mayor influencia en el rendimiento académico. Esta información servirá de apoyo para definir estrategias orientadas al fortalecimiento de la calidad educativa y al mejoramiento continuo de los procesos formativos.
Asimismo, la incorporación de técnicas estadísticas y modelos predictivos permitirá respaldar la toma de decisiones con información objetiva y basada en evidencia~\cite{tobon2013formacion}. Desde una perspectiva académica, el proyecto también aporta una aplicación práctica de la ingeniería de datos en el ámbito educativo, demostrando cómo el análisis de información puede convertirse en una herramienta valiosa para comprender mejor los resultados de aprendizaje y contribuir al fortalecimiento de la educación superior. Este sistema representa un avance significativo hacia la transformación digital en la gestión educativa, permitiendo a la universidad disponer de una plataforma sólida para el monitoreo, análisis y predicción del desempeño de sus estudiantes.

